\documentclass[10pt,a4paper]{amsart}
\usepackage[margin=19mm]{geometry}
\usepackage{amsmath,amssymb,mathtools}
\usepackage[scr=rsfs,cal=euler]{mathalfa}
\usepackage{xfrac}
\usepackage[unicode,hidelinks,bookmarks=false]{hyperref}
\newcommand{\ie}{i.e.\ }
\newcommand{\cf}{cf.\ }
\newcommand{\resp}{respectively\ }
\newcommand{\Subsection}{\subsection}
\providecommand{\nobreakdash}{\nobreak\hbox{-}}
\setlength{\emergencystretch}{2em}
\multlinegap0pt
\usepackage{tikz}
\usepackage{amssymb}
\usepackage{xcolor}
\usepackage{mathtools}
\newtheorem{lemma}{Lemma}
\def\ZE{\ensuremath{\mathbf E}}
\def\ZI{\ensuremath{\textbf 1}}
\def\ZR{\ensuremath{\mathbb R}}
\newcommand {\e }[1]{\eqref{#1}}
\title{On the best constants in Khintchine type \\inequalities for martingales}
\author{Grigori A. Karagulyan}
\date{Source: arXiv:2401.16153v1 (CC BY 4.0)}
\begin{document}
\maketitle

\noindent\emph{Source and licence.} On the best constants in Khintchine type \\inequalities for martingales. Version arXiv:2401.16153v1 (CC BY 4.0), distributed by arXiv under the Creative Commons Attribution 4.0 International licence, http://creativecommons.org/licenses/by/4.0/, as recorded in the arXiv licence metadata for this exact version and shown on the abstract page https://arxiv.org/abs/2401.16153v1. Full licence notice: \texttt{licenses/arxiv-2401.16153-CC-BY-4.0.txt}.\par\smallskip

\noindent\textit{Editorial scope.} Counted unit: the article's lemma L8 (source0.tex lines 511--517). The article's complete original proof of it is delivered with it (source0.tex lines 518--552, one \texttt{proof} environment of 35 lines). No mathematics is written, repaired or re-derived: the statement and the proof are the article's own lines, in the article's own notation, and no retained line is substituted or re-flowed.  No further source-local context is needed: every cross-reference of every retained line resolves inside the two delivered spans.  Nothing was omitted from any retained span, so the package carries no omission record.  No retained line contains a citation, so the wrapper carries no bibliography and no bibliography entry.  No retained line contains a figure environment, an image-inclusion command or drawing code, so no image is delivered and the package declares no asset dependency.  Editorial formatting only, in addition: an AMS \texttt{amsart} preamble that restates the article's own declarations and macros used by the retained lines, namely the environments \texttt{lemma} on separate counters, together with the standard packages those lines need.  The retained environments print in this document as Lemma 1 in order of appearance, whereas the article prints the same items with the numbering of its own full document.  The complete native source of the article as distributed in the arXiv e-print for this exact version is bundled unchanged as \texttt{sources/153674/source0.tex}.
\begin{lemma}\label{L8}
	Let $p\ge 1$and $\xi\in \ZR$ be fixed. Then for any non-trivial function $g\in L^1(0,1)$, satisfying $\ZE(g)=0$, the function
	\begin{equation}
		h(x)=\int_0^1\left|\xi+x\cdot g(t)\right|^pdt
	\end{equation}
is increasing over $(0,\infty)$.
\end{lemma}

\begin{proof}
Without loss of generality it is enough to prove that 
\begin{equation}\label{x52}
	\int_0^1\left|\xi+g(t)\right|^pdt<\int_0^1\left|\xi+r\cdot g(t)\right|^pdt\hbox{ for }r>1.
\end{equation}
Applying an approximation, we can suppose that $g$ is a step function on $[0,1)$ with equal constancy intervals. Let 
\begin{equation}
	E^+=\{x\in [0,1):\, g(x)\ge 0\},\quad E^-=\{x\in [0,1):\, g(x)<0\}.
\end{equation}
We will construct a sequence of step functions $g=g_1, g_2,\ldots, g_m=r\cdot g$, with the same constancy intervals as $g$ has, such that  
\begin{align}
	&\ZE(g_k)=0,\label{x50}\\
	&0\le g_k\cdot \ZI_{E^+}\le g_{k+1}\cdot \ZI_{E^+}\le rg(x)\cdot \ZI_{E^+},\label{x51}\\ 
	&0\ge  g_k\cdot \ZI_{E^-}\ge g_{k+1}\cdot \ZI_{E^-}\ge rg(x)\cdot \ZI_{E^-},\label{x48}\\
	&\int_X\left|\xi+g_k(t)\right|^pdt< \int_X\left|\xi+g_{k+1}(t)\right|^pdt,\quad k=1,2,\ldots, m-1.\label{x49}
\end{align}
Suppose by induction we have already defined the functions $g_1,\ldots, g_l$ such that the above conditions hold for $k=1,\ldots, l-1$. If $g_l=rg$, then the process will stop. Otherwise, applying \e{x50}, one can say that $g_l$ has at least two constancy intervals $I^+\subset E^+$ and $I^-\subset E^-$ such that
\begin{align}
	&rg(x)=b_+>g_l(x)=a_+\ge 0,\, x\in I^+,\\
	 &rg(x)=b_-<g_l(x)=a_-<0,\, x\in I^-.
\end{align}
Denote 
\begin{equation}
	\lambda=\min\{b_+- a_+, a_--b_-\}>0.
\end{equation}
We define $g_{l+1}$ by changing the values of $g_l$, on the intervals $I^+$ and $I^-$ by $a_++\lambda$ and $a_--\lambda$ respectively. Clearly, the conditions \e{x50}, \e{x51} and \e{x48} are satisfied for $k=l$ too. To show \e{x49} it is enough to observe the inequality
\begin{equation*}
	\int_{I^+\cup I^-}\left|\xi+g_l(t)\right|^pdt\le \int_{I^+\cup I^-}\left|\xi+g_{l+1}(t)\right|^pdt,
\end{equation*}
which is the same as the numerical inequality
\begin{equation*}
	|\xi+a_+|^p+|\xi+a_-|^p<|\xi+a_++\lambda|^p+|\xi+a_--\lambda|^p.
\end{equation*}
The latter follows from the fact that the function $t(x)=|c+x|^p+|c-x|^p$ is increasing when $x>0$. After this step of induction we will get one more interval (either $I^+$ or $I^-$), where $g_{l+1}$ coincides with $rg$. Thus, continuing the induction we will finally get a function $g_m=rg$. Thus, we will have \e{x52}, completing the proof of lemma.
\end{proof}

\end{document}
